\chapter{Заключение}

В работе были рассмотрены спектральные характеристики внутренних представлений трёх семейств моделей компьютерного зрения: ResNet, ViT и Swin. В фокусе исследования находился не только итоговый класс на выходе модели, а то, что происходит внутри: как по слоям меняется распределение энергии признаков по пространственным частотам. Иначе говоря, сохраняет ли модель более детальную, высокочастотную структуру или постепенно переводит представления в более сглаженный низкочастотный вид.

Для этого был собран единый протокол спектрального анализа. В случае ResNet анализировались карты признаков основных блоков. Для ViT последовательность токенов переводилась обратно в двумерную решётку патчей. Для Swin учитывалась его стадийная структура и изменение пространственного разрешения между уровнями. После получения представлений строился спектр мощности, затем выполнялось радиальное усреднение и нормализация. По полученным кривым рассчитывались три численные характеристики: спектральный центроид, доля низкочастотной энергии и скорость спектрального сжатия.

В начале исследования были сформулированы четыре гипотезы. Они и задавали логику экспериментальной части.

Первая гипотеза относилась к свёрточным сетям. Предполагалось, что ResNet по мере продвижения к глубоким блокам будет сдвигать спектр представлений в область низких частот. Эта гипотеза подтвердилась. Во всех рассмотренных моделях ResNet спектральный центроид снижается от ранних уровней к поздним, а доля низкочастотной энергии, наоборот, растёт. По числам это видно достаточно ясно: центроид меняется примерно от диапазона 0.17--0.21 на начальных блоках до 0.035--0.049 на глубоких. Доля низкочастотной энергии увеличивается примерно с 16--18\% до 30--31\%. Скорость сжатия у CNN также остаётся положительной и находится в диапазоне 0.46--0.51. По сути, ResNet ведёт себя как и ожидалось от иерархической свёрточной архитектуры: локальные детали на ранних уровнях постепенно уступают место более крупным и агрегированным признакам.

Вторая гипотеза касалась ViT. Я исходил из того, что Vision Transformer не обязан повторять поведение CNN, потому что он работает не со свёрточной пирамидой признаков, а с последовательностью патчевых токенов. Результаты это подтвердили. У ViT спектральный центроид не показывает устойчивого снижения. В отдельных моделях он, наоборот, растёт на средних или поздних блоках. Например, у ViT-large центроид увеличивается с 0.275 на начальном уровне до 0.437 на глубоком. С долей низкочастотной энергии картина тоже не такая ровная: у ViT-large она падает с 59.5\% до 30.6\%, а у ViT-base остаётся высокой даже на глубоких блоках. Значит, частотная динамика ViT устроена сложнее, чем простое движение к низким частотам. На неё, вероятно, одновременно влияют attention, MLP-блоки, нормализации и остаточные связи.

Третья гипотеза была связана со Swin Transformer. Предполагалось, что из-за patch merging и иерархической структуры Swin в конце будет приходить к более низкочастотным представлениям, но при этом его промежуточное поведение не обязано быть таким же прямым, как у CNN. Эта гипотеза подтвердилась, но с уточнением. На глубоких блоках Swin действительно даёт низкие значения центроида — примерно 0.048--0.054, а доля низкочастотной энергии выходит на уровень 29--30\%. Это близко к финальным значениям ResNet. Но на средних уровнях у части моделей центроид немного возрастает относительно начального уровня. Значит, Swin нельзя просто назвать ``CNN с вниманием''. Он скорее сочетает два механизма: иерархическое сжатие, похожее на CNN, и более сложное перераспределение признаков внутри трансформерных блоков.

Четвёртая гипотеза проверяла связь спектральных характеристик с точностью классификации. Здесь ожидалось, что прямой линейной зависимости не будет. Так и получилось. ResNet показывает самое выраженное спектральное сжатие и при этом достигает высокой точности в крупных вариантах. ViT сохраняет более широкий и местами немонотонный спектральный профиль, но это не приводит к автоматическому превосходству над лучшими ResNet и Swin. Swin, в свою очередь, хорошо работает на малых и средних масштабах, совмещая умеренное сжатие спектра с иерархической организацией признаков. Поэтому спектральные метрики полезны прежде всего как инструмент интерпретации. Они объясняют, как устроены представления, но не заменяют accuracy и не дают полного ответа о качестве модели.

Общий результат можно сформулировать так: архитектура заметно влияет на спектральную динамику внутренних представлений. ResNet даёт наиболее последовательное низкочастотное сжатие. ViT ведёт себя немонотонно: спектральная энергия может перераспределяться в сторону средних и более высоких частот даже на поздних блоках. Swin занимает промежуточную позицию: на финальных стадиях он становится близок к CNN по низкочастотному профилю, но по пути к этому состоянию демонстрирует более сложную динамику.

Практический смысл полученных результатов состоит в том, что предложенный протокол можно использовать для диагностики моделей компьютерного зрения. Он показывает, насколько быстро архитектура теряет или сохраняет частотное разнообразие признаков. Для задач, где мелкие локальные детали действительно несут полезную информацию, такой анализ может быть особенно полезен. Речь, например, о медицинских снимках, спутниковых изображениях, дефектоскопии, текстурах и мелких объектах.

У работы есть и ограничения. Спектральная энергия не равна информации в строгом смысле. Она не показывает напрямую, полезен ли конкретный частотный диапазон для классификации, не учитывает фазу спектра и не заменяет проверку качества на целевой задаче. Отдельное ограничение связано с ViT и Swin: перевод токенов в двумерную решётку нужен для сравнения с CNN, но такая решётка не является полной копией классической карты признаков. Токен уже содержит агрегированную информацию о патче, и это надо учитывать при интерпретации.

Дальше эту работу можно развивать в нескольких направлениях. Во-первых, полезно посмотреть, как спектральные метрики меняются не только у готовых моделей, но и в процессе обучения. Во-вторых, можно проверить связь спектрального профиля с устойчивостью к шуму и состязательным возмущениям. Ещё одно направление — перенести анализ на детекцию и сегментацию, где пространственная локализация признаков имеет большее значение, чем в обычной классификации. Наконец, на основе таких измерений можно пробовать строить регуляризаторы, которые будут контролировать сохранение или подавление отдельных частотных диапазонов.
